\documentclass[11pt]{article}
\usepackage[a4paper,margin=24mm]{geometry}
\usepackage{amsmath,amssymb,amsthm,microtype}
\usepackage[hidelinks]{hyperref}
\newcommand{\E}{\mathbb E}\newcommand{\Pp}{\mathbb P}
\newtheorem{theorem}{Theorem}\newtheorem{lemma}[theorem]{Lemma}
\title{A five-count reserve without a displacement error}
\author{}\date{28 October 2026}
\begin{document}\maketitle

The crossing identity gives nonnegative local losses. We bound five of
them directly, by replacing their kernels with a smooth function before
comparing conditional sums. Centering makes the replacement cost
proportional to the noise variance. There is no additive inward-displacement
error and no assumption that the two conditional laws coincide.

Let $U_0,U_1$ be centered with
\[
 -3/5\le U_0\le a_0,\qquad -a_1\le U_1\le3/5,\qquad
 a_0,a_1\ge0,\qquad a_0+a_1\le9/125.
\]
For $n\ge42$ and $p\in[2/5,53/100]$, write $q=1-p$,
$s_j=\E U_j^2$, $v=pq$ and $\lambda=n(qs_0+ps_1)\le9/25$.
All conditional sums below use independent coordinates. Set
\[
 h(x)=\E(x-U_0)_+-\E(x-1-U_1)_+-\Pp(1+U_1\le x),
\]
\[
 V_i=\sum_{l<i}U_{1l}+\sum_{l>i}U_{0l},\qquad
 g_i(t)=\E h(t-i+1-V_i).
\]
The retained crossing identity proves $g_i\ge0$ and
$\sum_i g_i=t-\E K$ whenever $na_0\le t<n(1-a_1)$.

\begin{theorem}[Five-count reserve]\label{thm:main}
Suppose $|t-np|<47/50$, and put $k=\lfloor t\rfloor$,
$a=1-5/n$ and $b=1+3/n$. Then
\begin{equation}\label{eq:reserve}
\sum_{i=k-1}^{k+3}g_i(t)\ge
\frac25\left\{\frac{1-e^{-3a\lambda}}2
-\frac{(9/25)b\lambda+3b^2\lambda^2}{16}
-\frac{10\pi^2\lambda}{n}\right\}.
\end{equation}
In particular the following uniform bounds hold:
\[
\begin{array}{c|rrrrrrr}
 n\text{ at least}&125&150&200&250&300&400&500\\ \hline
 \sum_{i=k-1}^{k+3}g_i/\lambda\text{ at least}
 &1/250&1/20&1/8&1/6&19/100&9/40&6/25 .
\end{array}
\]
At $\lambda=0$ these statements mean the corresponding nonnegative bound,
without division. The five indices lie in $\{1,\ldots,n\}$.
\end{theorem}

For the acquired raw thresholds $t=np+H\sqrt{v+\lambda/n}$ with
$|H|\le46/25$, the required displacement follows already at $n=42$:
\[
(46/25)^2(1/4+(9/25)/42)<(47/50)^2.
\]
The result therefore covers every phase and every threshold atom in the
right strip. It does not require the geometric budget or $d\le3/n$.
It is a lower bound for actual local losses, not a complete CDF comparison.
For any nonnegative weights $\pi_i$, it immediately yields a weighted
reserve by multiplying the right side of \eqref{eq:reserve} by
$\min_{k-1\le i\le k+3}\pi_i$. No distant-count tail is then needed.

\section{A compact minorant of each crossing kernel}
Define the continuously differentiable function
\[
B(x)=
\begin{cases}\sin^2(\pi x),&0\le x\le1,\\0,&\text{otherwise}.\end{cases}
\]
Its derivative is globally Lipschitz with constant $2\pi^2$:
the derivative vanishes at both joins and its derivative in each open
piece is bounded in absolute value by $2\pi^2$.

\begin{lemma}\label{lem:bump}
For every real $x$, one has $h(x)\ge(2/5)B(x)$.
\end{lemma}
\begin{proof}
For $x<1-a_1$, $h(x)=\E(x-U_0)_+\ge0$. For $x\ge1-a_1>a_0$,
put--call parity gives $h(x)=\Pp(1+U_1>x)-\E(1+U_1-x)_+$.
This is at least $(x-3/5)\Pp(1+U_1>x)\ge0$ until $x=8/5$,
and equals zero beyond that endpoint. Thus $h\ge0$ globally.
For $0\le x<1-a_1$, the terms involving $1+U_1$ vanish and Jensen gives
$h(x)=\E(x-U_0)_+\ge x$. We claim $B(x)\le(5/2)x$ on $[0,1]$.
For $x\le1/4$, use $\sin^2(\pi x)\le\pi^2x^2$ and $\pi^2<10$.
For $x\ge2/5$, use $B\le1$. Between these points, concavity and the
tangent at $1/4$ give
\[
B(x)\le\tfrac12+\pi(x-\tfrac14)\le\tfrac52x.
\]
The last affine inequality holds at both endpoints, using $\pi<10/3$.

For $1-a_1\le x\le1$, put $D=1+U_1$ and $L=\E(D-x)_+$.
Put--call parity gives $h(x)=\Pp(D>x)-L$, with a strict survival event.
Since $D-x\le1+3/5-x$ on this event and $\E D=1$,
\[
h(x)\ge\frac{x-3/5}{1+3/5-x}L
\ge\frac{41}{84}(1-x).
\]
Here $x\ge1-a_1\ge116/125$ and $L\ge1-x$.
On the other hand
\[
\frac25B(x)\le\frac25\pi^2(1-x)^2
\le\frac25\left(\frac{22}{7}\right)^2\frac9{125}(1-x)
\le\frac{41}{84}(1-x).
\]
At $x=1$ the minorant is zero. This includes all atoms of $D$ because
the displayed survival event stays strict.
\end{proof}

\section{Four centered replacements suffice}
We use the following elementary smooth comparison. If $F'$ is Lipschitz
with constant $L$ and $Z$ is centered, then
\[
|\E F(r+Z)-F(r)|\le \tfrac L2\E Z^2.
\]
Taylor's integral remainder proves this even if $F''$ exists only
piecewise. Conditional on an independent random $r$, it follows that
replacing a centered summand of variance $s_0$ by one of variance $s_1$
changes the expectation by at most $L(s_0+s_1)/2$.

Write $u=t-k\in[0,1)$ and freeze the middle omitted-coordinate sum
\[
V=V_{k+1}\ \stackrel{\mathrm{law}}=
U_{11}+\cdots+U_{1k}
+U_{0,k+2}+\cdots+U_{0n}.
\]
For $j=-2,-1,0,1,2$, let $V^{(j)}$ have the law of $V_{k+1+j}$,
and put $B_j(x)=B(u-j-x)$. All five sums have $n-1$ terms.
For $j>0$, the law of $V^{(j)}$ replaces $j$ copies of $U_0$ by $U_1$.
The two expectation differences regroup exactly as
\begin{align*}
&\E B_1(V^{(1)})+\E B_2(V^{(2)})
 -\E B_1(V)-\E B_2(V)\\
&=\{\E(B_1+B_2)(V^{(1)})-\E(B_1+B_2)(V)\}\\
&\hspace{12mm}+\{\E B_2(V^{(2)})-\E B_2(V^{(1)})\}.
\end{align*}
Each brace is one centered replacement. The two translated bumps in
$B_1+B_2$ have disjoint open supports, so its derivative still has
Lipschitz constant $2\pi^2$, rather than twice that value.
The negative indices have the identical regrouping, now replacing $U_1$
by $U_0$. The derivatives are zero at every common boundary.
Consequently
\begin{equation}\label{eq:swap}
\sum_{j=-2}^2\E B(u-j-V^{(j)})
\ge\sum_{j=-2}^2\E B(u-j-V)-4\pi^2(s_0+s_1).
\end{equation}
There are enough copies of each law: $k-1\ge1$ and $k+3\le n$ follow
from $n\ge42$, the specified $p$ range and the raw displacement bound.
No conditional sum at one count is assumed independent of a sum at another.

The translates partition a periodic sine square:
\[
\sum_{j\in\mathbb Z}B(u-j-z)=\sin^2(\pi(u-z)).
\]
Every omitted nonzero term, with $|j|\ge3$, requires $|z|>2$.
At integer joins both adjacent bumps vanish. Therefore
\begin{equation}\label{eq:tail}
\sum_{j=-2}^2\E B(u-j-V)
\ge \E\sin^2(\pi(u-V))-\frac{\E V^4}{16}.
\end{equation}

\section{Variance and a single characteristic function}
Put $\delta=k-np\in(-97/50,47/50)$. The actual frozen variance is
\[
\tau=\E V^2=(nq-\delta-1)s_0+(np+\delta)s_1,
\qquad a\lambda\le\tau\le b\lambda.
\]
Indeed $p\ge2/5$, $q\ge47/100$, so the lower coefficient losses are
at most $5/n$ and the upper gains at most $3/n$. Since $|U_j|\le3/5$,
independence gives
\[
\E V^4\le(9/25)\tau+3\tau^2
\le(9/25)b\lambda+3b^2\lambda^2.
\]
Also $s_0+s_1\le5\lambda/(2n)$.

For completeness, for $|x|\le84/125$,
\[
1-\cos(2\pi x)\ge3x^2.
\]
On $0\le x\le2/3$, concavity gives
$\sin(\pi x)\ge(3\sqrt3/4)x\ge(51/40)x$.
On $2/3\le x\le84/125$, the Lipschitz bound for sine gives
$\sin(\pi x)>17/20-44/2625$.
Squaring these positive estimates proves the claim by rational
arithmetic, using $\sqrt3>17/10$ and $\pi<22/7$.

For either conditional noise $Z$ and its independent copy $Z'$,
$|Z-Z'|\le3/5+9/125=84/125$. It follows that
\[
1-|\E e^{2\pi iZ}|^2
=\E\{1-\cos(2\pi(Z-Z'))\}\ge6\operatorname{Var}Z.
\]
Thus $|\E e^{2\pi iZ}|\le e^{-3\operatorname{Var}Z}$. Independence
of the $n-1$ frozen summands, not a claim about the selected count, gives
\[
\E\sin^2(\pi(u-V))
\ge\frac{1-|\E e^{2\pi iV}|}{2}
\ge\frac{1-e^{-3\tau}}2
\ge\frac{1-e^{-3a\lambda}}2.
\]
Combine this with Lemma~\ref{lem:bump}, \eqref{eq:swap}
and \eqref{eq:tail} to prove \eqref{eq:reserve}.

To verify the table without floating-point exponentials, put $L=9/25$,
$a_N=1-5/N$, $b_N=1+3/N$ and
$P_6(x)=\sum_{j=0}^6x^j/j!$.
Concavity of $1-e^{-3a_N\lambda}$ on $[0,L]$ gives the rational lower
coefficient
\[
c_N=\frac25\left\{
\frac{1-P_6(3a_NL)^{-1}}{2L}
-\frac{(9/25)b_N+3b_N^2L}{16}
-\frac{10(22/7)^2}{N}\right\}.
\]
Here $e^x\ge P_6(x)>0$, and increasing $n$ improves each term.
Each $c_N$ strictly exceeds its displayed table value by exact rational
arithmetic. The argument is uniform down to zero noise.

\section*{Scope and dependencies}
Only elementary inequalities and the local crossing kernel are used.
The kernel's nonnegativity is also proved above by the same separated-support
argument as in the retained v13 crossing note. Identifying the full sum
with $t-\E K$, or joining these local losses to the actual CDF, uses that
note's exact identities with their endpoint conventions.
The five-count reserve needs no external probability theorem, finite
certificate, acquisition argument, common conditional law, or selected-index
independence. It does not itself settle the remaining normal, curvature and
binomial-weight corrections. The table does not claim a positive bound
for every $42\le n<125$.
\end{document}

% Deliberate reviewed-source mutation.
